\documentclass[11pt]{article}
\usepackage[margin=1in]{geometry}
\usepackage{listings}
\usepackage{xcolor}
\usepackage{booktabs}
\usepackage{hyperref}

\hypersetup{colorlinks=true, linkcolor=blue, urlcolor=blue}

\lstset{
  basicstyle=\ttfamily\small,
  keywordstyle=\color{blue},
  commentstyle=\color{gray},
  frame=single,
  breaklines=true,
  columns=fullflexible
}

\title{FPGA Session Log --- \texttt{final\_8}: Hardware Bring-Up \& Debug\\
\large Digilent Arty A7-35T \,\textperiodcentered\, Open-Source Toolchain}
\author{Lucas Zhou}
\date{July 16--17, 2026}

\begin{document}
\maketitle

\section{Session Goal}
Build, flash, and verify the finished \texttt{top.v} (2-bit $\times$ 2-bit
multiplier shown on a Pmod SSD seven-segment display) against the 2021
ECE~15A final, problem~8 --- and debug why the display showed wrong or
partial digits on real hardware.

\section{Simulation (passed, unchanged all session)}
The Icarus Verilog sweep of all 16 input combinations matches the case
table in \texttt{top.v} exactly: products 0--9 as expected
($A\times B$ for $A,B \in \{0..3\}$). The logic was correct from the
start; every bug below was a pin-mapping issue in \texttt{arty.xdc}.

\begin{lstlisting}[language=bash]
source tools/oss-cad-suite/environment
iverilog -o tb_top.vvp tb_top.v top.v
vvp tb_top.vvp        # 16/16 rows match the case table
\end{lstlisting}

\section{Debug Timeline}
\begin{enumerate}
  \item \textbf{Flash \#1 (XDC v1, seg[N] $\to$ AA..AG alphabetical):}
    scrambled shapes (a ``7''-like glyph, fragments). Cause:
    \texttt{top.v} numbers segments \emph{geometrically bottom-to-top}
    (seg[0..6] = D,C,E,G,B,F,A), not alphabetically.
  \item \textbf{Flash \#2 (XDC v2, permuted, still JA-only):} input
    \texttt{1011} ($2\times3$) showed only \emph{part} of a 6.
  \item \textbf{Photo analysis (\texttt{IMG\_8616}):} two findings ---
    (a) the Pmod SSD straddles \emph{two} connectors: header J1
    (AA--AD) in JA's top row, header J2 (AE, AF, AG, CAT) in
    \emph{JB's} top row, so v1/v2 drove three segments + digit select
    on untouched JA bottom-row pins (they floated); (b) the glyphs are
    mounted 180$^\circ$ rotated (decimal points top-left), so each
    signal lights the segment opposite its name.
  \item \textbf{Flash \#3 (XDC v3, current):} AE/AF/AG/CAT moved to JB
    pins E15/E16/D15/C15; all seven assignments rotated 180$^\circ$.
    Digit select \texttt{C} now actually driven by
    \texttt{assign C = 1'b1}.
\end{enumerate}

\section{Final Pin Map (XDC v3)}
\begin{center}
\begin{tabular}{@{}lllll@{}}
\toprule
Port & Meaning & SSD signal & Pin & Connector \\
\midrule
\texttt{A2/A1/B2/B1} & inputs & --- & A10/C10/C11/A8 & SW3/SW2/SW1/SW0 \\
\texttt{seg[0]} & bottom       & AA & G13 & JA1 \\
\texttt{seg[1]} & lower-right  & AF & E16 & JB2 \\
\texttt{seg[2]} & lower-left   & AB & B11 & JA2 \\
\texttt{seg[3]} & middle       & AG & D15 & JB3 \\
\texttt{seg[4]} & upper-right  & AE & E15 & JB1 \\
\texttt{seg[5]} & upper-left   & AC & A11 & JA3 \\
\texttt{seg[6]} & top          & AD & D12 & JA4 \\
\texttt{C} & digit select      & CAT & C15 & JB4 \\
\bottomrule
\end{tabular}
\end{center}

\section{Open Items}
\begin{itemize}
  \item \textbf{Visual confirmation} of XDC v3 on hardware (expect a
    complete upright 6 for input \texttt{1011}; spot-check 9, 4, 1, 0).
  \item If the \emph{other} digit of the two should light: flip
    \texttt{assign C = 1'b1} to \texttt{1'b0} in \texttt{top.v},
    rebuild.
  \item Current load is \textbf{SRAM only} (lost on power-off). For the
    final demo, write to flash:
    \texttt{openFPGALoader -b arty -f top.bit}.
\end{itemize}

\section{Files in \texttt{projects/final\_8/}}
\begin{tabular}{@{}lll@{}}
\toprule
File & Status & Notes \\
\midrule
\texttt{top.v} & unchanged, correct & author's multiplier + case table \\
\texttt{arty.xdc} & rewritten twice (v3) & JA/JB split + 180$^\circ$ rotation \\
\texttt{xdc-explained.pdf} & new this session & line-by-line XDC tutorial + full debug story \\
\texttt{session-log-2026-07-13.pdf} & prior session & syntax fixes, testbench, sim flow \\
\texttt{top.bit} & flashed to SRAM & 3 builds, all 0 errors \\
\bottomrule
\end{tabular}

\vspace{1em}
\noindent For the beginner-level, line-by-line explanation of the XDC
syntax and all three versions, see \texttt{xdc-explained.pdf}.

\end{document}
